% ============================================================
\section{Task Formulation and the Evidence Trace}
\label{sec:method}

An example consists of a query $q$ and a retrieved set $D=\{d_1,\dots,d_N\}$. The
system must produce either a citation-linked answer or the controlled refusal
string \texttt{CANNOT ANSWER, INSUFFICIENT EVIDENCE}. It reasons only over the
retrieved evidence and refuses when that evidence is insufficient. Rather than
mapping $(q,D)$ directly to a response, we supervise an intermediate
representation emitted as part of the output.

\paragraph{Stage 1: document adjudication.} Every $d_i$ receives a verdict in
\{\textit{supports}, \textit{partially supports}, \textit{irrelevant}\}, a
one-sentence key fact entailed by the snippet, and a coarse source-quality label.
The \textit{partially supports} class is deliberately broad, covering on-topic
but incomplete, hedged, scoped, or single-facet evidence. This choice avoids
the failure mode we care about: a binary useful/useless split causes refusal
whenever no single snippet is individually decisive.

\paragraph{Stage 2: conflict analysis.} The Stage-1 judgments are aggregated into
exactly one of five labels, each prescribing a different aggregation strategy.
Table~\ref{tab:taxonomy} (Appendix~\ref{app:cats}) gives each type's definition
and the response behaviour it requires. The taxonomy separates
\emph{complementarity} from \emph{conflict}: two snippets can contribute
different necessary pieces without disagreeing.

\paragraph{Stage 3: grounded synthesis.} Stages 1 and 2 are emitted \emph{inside}
the \verb|<think>| block. Stage 3 is the response itself and sits outside it. The
final response is conditioned on the conflict label, with sentence-level
\texttt{[dX]} citations, or is the exact refusal string, terminating with
\verb|[[END-OF-ANSWER]]|. Figure~\ref{fig:system} places this in the overall
pipeline.

\section{Dataset Generation and Validation}
\label{sec:data}

Table~\ref{tab:data} summarizes the generated dataset and its validation.
Appendix~\ref{app:data} gives the complete construction record, per-source
distributions, and retrieval protocol; Appendix~\ref{app:human-review} gives
the human-review statistics.

\begin{table}[htbp]
\centering
\begin{minipage}[t]{0.47\linewidth}\vspace{0pt}
\centering
{\tabstyle\setlength{\tabcolsep}{3pt}
\begin{tabular*}{\linewidth}[t]{@{\extracolsep{\fill}}l rr rr@{}}\toprule
& \multicolumn{2}{c}{\textbf{Train}} & \multicolumn{2}{c}{\textbf{Benchmark}}\\
\cmidrule(lr){2-3}\cmidrule(lr){4-5}
\textbf{Conflict type} & $n$ & \% & $n$ & \%\\\midrule
No conflict & 343 & 36.4 & 211 & 28.7\\
Complementary & 286 & 30.3 & 221 & 30.0\\
Conflicting opinions & 158 & 16.8 & 109 & 14.8\\
Outdated & 136 & 14.4 & 158 & 21.5\\
Misinformation & 20 & 2.1 & 37 & 5.0\\
\midrule
Answerable & 649 & 68.8 & 608 & 82.6\\
Refusal-required & 294 & 31.2 & 128 & 17.4\\
\midrule
\textbf{Records} & \textbf{943} & & \textbf{736} & \\
Retrieved documents & 6,642 & & 3,701 & \\
Docs / record (mean) & 7.04 & & 5.03 & \\
Unique source domains & 2,633 & & 1,431 & \\
\bottomrule\end{tabular*}}

\end{minipage}%
\hspace{0.03\linewidth}%
\begin{minipage}[t]{0.47\linewidth}\vspace{0pt}
\centering
{\tabstyle\setlength{\tabcolsep}{0pt}
\begin{tabular}[t]{@{}p{\linewidth}@{}}\toprule
\textbf{Grounded Refusal (GR)}\newline whether the system answered exactly when the evidence permits answering; scored deterministically from the answerability label\tabularnewline
\addlinespace[9.2pt]
\textbf{Behaviour Adherence (BA)}\newline whether the answer followed the response policy its conflict type requires\tabularnewline
\addlinespace[9.2pt]
\textbf{Factual Grounding (FG)}\newline whether each extracted claim is both supported by an eligible document \emph{and} cited\tabularnewline
\addlinespace[9.2pt]
\textbf{Single-Truth Recall (STR)}\newline whether the response adopts the target answer as its own conclusion, rather than quoting or attributing it\\
\bottomrule\end{tabular}}

\end{minipage}
\dualcaption{table}{tab:data}{tab:cats}%
{\textbf{Dataset composition and scale.} Conflict-type and answerability distribution of the training split (943 records, including validation) and the held-out benchmark (736). The benchmark is deliberately harder than the training distribution. It has fewer no-conflict cases ($28.7\%$ against $36.4\%$), a $1.5\times$ larger outdated-information share, and a $2.4\times$ larger misinformation share. Full per-source composition in Appendix~\ref{app:data}}%
{\textbf{The four CATS components}, each reported over its own applicability denominator because the constructs apply to different subsets of examples. Appendix~\ref{app:cats} gives the full definitions, FG eligibility rules, per-type BA policies and the committee voting rule.}
\end{table}

\subsection{Training and validation}

The 943-example training release (862 train / 81 validation) combines three
source families: \textsc{Conflicts} \citep{cattan2025dragged} and held-out
refusal examples derived from Trust-Align \citep{song2025trustalign}. The
third family draws on examples from the benchmark construction stream that
were \emph{not} admitted to the final benchmark. We checked these examples for
identifier overlap against the released benchmark and found the intersection
empty. Appendix~\ref{app:data} gives the per-family counts.

\subsection{An independent benchmark}

The 736-example benchmark is built in three tracks and reporting them
separately is deliberate. \textbf{(1) Track A: web-retrieved (508).} Queries are drawn from ConflictingQA
\citep{wan2024evidencelanguagemodelsconvincing}, SituatedQA
\citep{zhang2021situatedqaincorporatingextralinguisticcontexts}, FreshQA
\citep{vu2024freshllms} and QACC \citep{liu2025opendomainquestionanswering}, with
every query overlapping the training pool excluded. Documents are retrieved live
and a query-focused window is selected by a TAS-B dense bi-encoder
\citep{hofstatter2021tasb}. Appendix~\ref{app:data} gives the full funnel and the
five-document sampling rule. \textbf{(2) Track B: refusal (128).} Refusal-required examples come from the
held-out refusal pool of Trust-Align \citep{song2025trustalign} and are
\emph{not} web-retrieved. \textbf{(3) Track C: curated conflict (100).} Open-web
retrieval cannot reliably surface every conflict regime, such as a fact whose
truth value has since changed. We therefore add 78 examples built from
dated Wikipedia revision pairs, in the spirit of \citet{hou2024wikicontradict},
19 multi-hop contexts from HotpotQA \citep{yang2018hotpotqa} and three
medical/misinformation cases. The track is strongly class-targeted, cleanly
separating outdated-information and complementary cases
(Table~\ref{tab:tracks}, Appendix~\ref{app:data}).

\subsection{Annotation Depth \& Human Review}

Each of the $10{,}343$ retrieved documents carries six fields: an adjudicated
\textit{verdict}, an extracted \textit{key fact}, and a snippet-anchored
\textit{supporting quote}. It also carries a \textit{rationale}, a
\textit{source-quality label}, and the complete \textit{annotator vote tally}.
Labels come from a weighted multi-annotator committee that votes on the
categorical decision at each stage. The committee then adopts the winning
side's explanation bundle intact from its highest-weight member, keeping each
verdict, quote, and rationale internally coherent. The annotation contract
requires the key fact to be entailed by its quote. It also requires the quote
to be a contiguous verbatim span of at most fifty words. Appendix~\ref{app:annotation}
gives the formal aggregation rule and the full contract. The prompts that
enforce it are released with the resource (Appendix~\ref{app:annprompts}).

Two human workflows, kept separate, cover the two populations. Seven
reviewers screened the benchmark candidates on a retention decision, a
preliminary conflict type, and eight quality judgments
(Appendix~\ref{app:human-review} lists each). This screening reached
$94.77\%$ five-way conflict-type agreement (Cohen's $\kappa=0.92$) and $98.21\%$
retention agreement ($\kappa=0.92$). Because this second pass is
\emph{sequential}, the figures measure verification consistency rather than
blind double annotation. Training-side review (two reviewers per example)
reached $83.46\%$ conflict-type agreement ($\kappa=0.77$). A separate,
class-balanced 100-pair study checks the Stage-1 per-document verdict itself,
not the downstream conflict-type label. Three reviewers agree at Fleiss'
$\kappa=0.83$ and agree with the committee at $92\%$ ($\kappa=0.88$ against
the human majority). Appendix~\ref{app:human-review} and
Appendix~\ref{app:stage1review} give the full protocols.

% ============================================================
\section{Trace-Supervised Fine-Tuning}
\label{sec:sft}

\subsection{Objective and prompt variants}

The supervision target is the serialised trace followed by the final response.
Three prompt variants are used: the \textbf{strict} prompt states the full
taxonomy, verdict heuristics, and output contract ($1{,}965$ words). The
\textbf{runtime} prompt keeps a compact structural schema ($195$ words). The
\textbf{minimal} prompt supplies only the query, the documents, and a
four-sentence evidence policy ($40$ words). The mixture also supervises
document-verdict, conflict-type, and final-response prediction as separate
task views. This gives each component explicit signal in isolation as well as
inside the full trace. We train and evaluate across all three variants to
test whether conflict-aware behaviour is \emph{learned} or merely
\emph{recited} from the instruction (\S\ref{sec:internalization}).
Appendix~\ref{app:e2eprompts} gives the three prompts verbatim and their
design rationale.

\subsection{Fine-tuning}
All models are fine-tuned with QLoRA
\citep{dettmers2023qloraefficientfinetuningquantized,hu2021loralowrankadaptationlarge},
a parameter-efficient method that trains low-rank adapters (LoRA) atop a
frozen, quantized base model. Training uses the standard
token-level cross-entropy of the trace and response, with no auxiliary loss.
Appendix~\ref{app:sft} gives the hyperparameters, the checkpoint selection
criterion, and the two training mixtures (recipes~K and~L). These mixtures
differ by model family and are needed to interpret the per-model results in
\S\ref{sec:results}.

% ============================================================
\section{CATS: Conflict-Aware Trust Score}
\label{sec:cats}

Standard RAG metrics collapse distinct failures. Answering when evidence is
absent, mishandling the conflict type, making unsupported claims, and omitting
the target answer are all ``a bad answer''. CATS separates them into four
components with independent applicability.

Table~\ref{tab:cats} states what each component measures and the population it is
scored over. The definitions are operational. FG counts a claim only when the
supporting document's \emph{gold} verdict is \textit{supports} or \textit{partially
supports}. STR excludes conflicting-opinion items because no single conclusion
exists for them. Appendix~\ref{app:cats} sets out each rule. The denominators
differ, so any comparison against another framework should start there.

\textbf{Judges:} BA, FG, and STR are judged by a committee of three locally
served open-weight models with fixed weights, aggregated by weighted majority.
Factual-grounding support additionally requires corroboration by at least two
judges, and GR is never committee-judged. We compare this against human review
in \S\ref{sec:human}. The complete judge bundle is one system message plus
three task prompts and the five-type behaviour rubric, reproduced verbatim in
Appendix~\ref{app:catsprompts}. Metric definitions and aggregation are in
Appendix~\ref{app:cats}. \textbf{Aggregate scores:} the four components can also be combined into a
single number per example, gated by the answer/refuse decision so that a wrong
decision scores zero. We define two averages, \textbf{CATS-Prevalence} and
\textbf{CATS-Balanced}, in Appendix~\ref{app:aggregate}. We do not report
either in the main results. Because a refusal-required example is scored on
its decision alone, such examples carry $40\%$ of CATS-Balanced while making
up $17.4\%$ of the benchmark. A model can therefore raise the aggregate by
improving refusals without improving any answer. Every claim in
\S\ref{sec:results} is stated on the components.

% ============================================================
\section{Experimental Setup}
\label{sec:setup}

\subsection{Primary evaluation}

We fine-tune four instruction-tuned open-weight models, Qwen2.5-7B and
Qwen2.5-32B \citep{yang2024qwen25}, Llama-3.1-8B \citep{dubey2024llama3} and
Mistral-7B \citep{jiang2023mistral7b}, and compare each against its baseline
version on the four CATS components. Within a cell, the benchmark examples,
prompts, decoding policy, retry policy, post-processing, and evaluator are
identical. Only the LoRA adapter differs (Table~\ref{tab:main}). Each pair is
evaluated under the strict and minimal prompts, with the intermediate runtime
prompt reported in Appendix~\ref{app:full}. Generation is deterministic
(temperature $0$) with profile-specific continuation budgets and
sentinel-based stopping. Retries only enlarge the token budget for the
identical prompt, never resampling or repairing content
(Appendix~\ref{app:inference}). Per-example scores lie in $[0,1]$, so a
distribution-free 95\% interval half-width at $n=736$ is $\pm0.034$. We treat
differences below this threshold as uninformative. We report exact sign
tests for paired comparisons across the twelve model$\times$prompt cells, and
use exact Clopper--Pearson (binomial) intervals on the 128-example refusal
subgroup.

\subsection{Ablation studies}

\textbf{Oracle ablations.} To test whether the three stages are separable, we
re-run the fine-tuned models with gold information injected one stage at a
time (\S\ref{sec:oracle}). The \textbf{+conflict label} condition supplies the
gold Stage-2 type, and \textbf{+per-doc notes} supplies the gold Stage-1
verdicts and key facts. The \textbf{+both} condition supplies each.
\textbf{Comparison with prompting methods.} \emph{Chain-of-Note}
\citep{yu2024chainofnote} writes per-document reading notes before answering,
the closest analogue to our Stage~1. \emph{Few-shot Chain-of-Thought}
\citep{wei2022cot} supplies worked exemplars of the reasoning format. Both
receive the identical benchmark, decoding policy, and evaluator
(\S\ref{sec:comparators}).
